\section{Summary: strengths and limitations of \paper{}}\label{sec:summary}

\subsection*{What is sound}
\begin{itemize}[leftmargin=2em]
\item \emph{The framework} (Section~\ref{sec:extremal}). Extremal series, the defect, the
inequality $\abs{\ip{T^jv}{x_0-v}}\le\tau$, and the criterion that a norm limit of extremal
images with defect tending to $0$ is non-cyclic. This is a clean and, as far as we know, new
reformulation: the problem becomes to find a norm-Cauchy sequence of approximate solutions.
\item \emph{The confinement identity} (Lemma~\ref{lem:confinement}). Holding one inner product
fixed does control the norm distance between approximate solutions. This is the best idea in
the paper and it is correct.
\item \emph{The type 1 / type 2 dichotomy} and the reductions.
\item \emph{Dominance of the leading coefficient} (Lemmas 2--5; Section~\ref{sec:domination})
and the \emph{Case I} tail bound (Section~\ref{sec:scaling}), after corrections.
\item \emph{The first-order variation formulas} (Section~\ref{sec:variation}), after
corrections.
\item \emph{The type 2 reduction} (Section~\ref{sec:type2}): for a fixed $x_0$ it produces
extremal images at distance about $0.3$ from $x_0$ with defect tending to $0$, and the
weak-limit alternative of the endgame is valid when it applies.
\end{itemize}

\subsection*{What is wrong but repairable}
Lemma~1 (constants), the bookkeeping in Lemma~2, the explicit constants in (41)--(42), the
direction of the argument for (44), and several formula slips; see the table in
Section~\ref{sec:gaps}. Lemma~1 is in fact unnecessary (Remark~\ref{rem:trivialstart}).

\subsection*{What is wrong and not repaired}
\begin{itemize}[leftmargin=2em]
\item \emph{Case II} ((26)--(27)): the estimate overlooks a first-order term
(Remark~\ref{rem:caseIIerror}), and (27) is false for type 1 operators at small defect
(Proposition~\ref{prop:27false}). Whether the defect still decreases under the scaling
iteration is open.
\item \emph{The Main Construction} ((31)--(45)): only first-order; no step size, no proof that
it continues, and a requirement of relative precision $\tau^{15}$ (Gap~\ref{gap:secondorder});
real versus complex independence (Gap~\ref{gap:real}).
\item \emph{The convergence step} ((46) and after): the justification is refuted.
Theorem~\ref{thm:abstract-nogo} (machine-checked) and Corollary~\ref{cor:explicit-nogo} show
that in an operator satisfying every hypothesis the argument uses, there are admissible states
of arbitrarily small defect from which the confined construction cannot continue. (For the states of these notes the non-degeneracy hypotheses of \paper{} are verified only
numerically, Remark~\ref{rem:33}; the standalone paper \texttt{tex/refutation.tex} gives
explicit states for which they are proved, and shows these are reachable from the starting
vector of \paper{}.) The type 2 endgame depends on this step too (Gap~\ref{gap:type2}).
\end{itemize}

\subsection*{Why we think the method cannot be repaired as it stands}
\begin{enumerate}
\item \emph{Small defect is cheap and does not localise.} A small defect follows from
$0\in\sigma_{\mathrm{ap}}(T^*)$ alone (Remark~\ref{rem:trivialstart}), and in the model there
are approximate solutions of arbitrarily small defect far from every solution and far from each
other (Corollary~\ref{cor:obstruction}).
\item \emph{The hypotheses of the paper do not see the difference.} Type 1 constants,
extremality, dominance and the size of the defect are the same for states that can be
completed and states that cannot. The only quantity we know that does see it is $D_T$ of
Section~\ref{ssec:abstract}, which is a global object.
\item \emph{The natural construction drifts the wrong way.} Started next to a true solution in
the model, least-norm steps lower the defect by oscillation and move away from the solution
(Section~\ref{sec:experiments}).
\item \emph{The confinement idea needs a solution to exist nearby.} By
Lemma~\ref{lem:constraintconfines} the constrained sequence stays in a small ball, so it can
succeed only if that ball contains a vector of defect $0$. In a proof by contradiction there is
none, so the whole weight rests on proving that constrained steps always exist, and that is
exactly the statement the model refutes for the reasons given.
\end{enumerate}
A proof along these lines would need a new ingredient: a rule for choosing steps and a quantity,
monotone or at least controlled along the resulting trajectory, that certifies continuation.
Nothing of the kind is in \paper{}. Remark~\ref{rem:where} records the limit of what we have
shown: the obstruction is exhibited in operators with many reducing subspaces, and we have no
information about what the construction would do for a transitive operator.

\subsection*{What is machine-checked}
In the Lean project at the repository root (Mathlib, no \texttt{sorry}): the non-cyclicity
criterion; the confinement identity; the pointwise identities; the abstract rigidity
inequality and its self-adjoint form; the abstract no-continuation theorem. Not formalised:
Sections~\ref{sec:extremal}--\ref{sec:domination}, the instantiation of the no-continuation
theorem in the model, and everything numerical. All hand proofs were re-derived once by an
independent referee pass.
